\documentclass[a4paper,12pt,twoside,openright,titlepage]{article}

%Additional packages
\usepackage[utf8]{inputenc}
\usepackage[T1]{fontenc}
\usepackage[dutch,english]{babel}
\usepackage{syntonly}
\usepackage[official]{eurosym}

% Handle images
%\usepackage[graphicx]
\usepackage{graphicx}
\graphicspath{ {./images/}{./styles/} }
\usepackage{float}
\usepackage{wrapfig}

% Handle URLs
\usepackage{xurl}
\usepackage{hyperref}
\hypersetup{colorlinks=true, linkcolor=blue, citecolor=blue, filecolor=blue, urlcolor=blue, pdftitle=, pdfauthor=, pdfsubject=, pdfkeywords=}

% Tables and listings
\usepackage{multirow,tabularx}
\usepackage[table]{xcolor} % Table colors
\usepackage{scrextend}
\addtokomafont{labelinglabel}{\sffamily}
\usepackage{listings}
\usepackage{adjustbox}

% Turn on indexing
\usepackage{imakeidx}
\makeindex[intoc]

% Define colors
\usepackage{color}
\definecolor{ashgrey}{rgb}{0.7, 0.75, 0.71}



% Listing style
\lstset{
  backgroundcolor=\color{ashgrey}, % choose the background color; you must add \usepackage{color} or \usepackage{xcolor}; should come as last argument
  rulecolor=\color{black},         % if not set, the frame-color may be changed on line-breaks within not-black text (e.g. comments (green here))
  frame=single,	                   % adds a frame around the code
  basicstyle=\ttfamily,  % the size of the fonts that are used for the code
  extendedchars=true,    % lets you use non-ASCII characters; for 8-bits encodings only, does not work with UTF-8
  breakatwhitespace=true, % sets if automatic breaks should only happen at whitespace
  breaklines=true,        % sets automatic line breaking
  keepspaces=true,        % keeps spaces in text, useful for keeping indentation of code (possibly needs columns=flexible)
  columns=fullflexible,	  % Make sure no extra spaces are added.
  showstringspaces=false, % if true show spaces in strings adding particular underscores
  showspaces=false        % show spaces everywhere adding particular underscores; it overrides 'showstringspaces'
}



% Uncomment for production
% \syntaxonly

% Style
\pagestyle{headings}

% Turn on indexing
\makeindex[intoc]

% Define document
\author{D. Leeuw}
\title{Linux: Werken met bestanden op de CLI}
\date{\today\\v.0.9.5\\
\vfill
\raggedright
\copyright\ 2020-2025 Dennis Leeuw\\

%\begin{figure}
\includegraphics[width=0.3\textwidth]{CC-BY-SA-NC.png}
%\end{figure}

Dit werk is uitgegeven onder de Creative Commons BY-NC-SA Licentie en laat anderen toe het werk te kopi\"eren, distribueren, vertonen, op te voeren, en om afgeleid materiaal te maken, zolang de auteurs en uitgever worden vermeld als maker van het werk, het werk niet commercieel gebruikt wordt en afgeleide werken onder identieke voorwaarden worden verspreid.


}


\begin{document}
\selectlanguage{dutch}

\maketitle

%%%%%%%%%%%%%%%%%%%
%%% Introductie %%%
%%%%%%%%%%%%%%%%%%%

\section{Over dit Document}
\subsection{Leerdoelen}
Na het bestuderen van dit document heeft de lezer kennis van:
\begin{itemize}
\item hoe er gewerkt kan worden met bestand en directories
\item de commando's: touch, cat, mkdir, rmdir, rm
\item het werken met redirects (onleidingen): >, >>, <, <<
\item de betekenis van stdin, stdout en stderr
\end{itemize}

Dit document sluit aan op de volgende onderdelen van de LPI:
\begin{itemize}
\item LPI Linux Essentials 010-160 - 6.2.3 Using Directories and Listing Files (weight: 2)
\item LPI Linux Essentials 010-160 - 6.2.4 Creating, Moving and Deleting Files (weight: 2)
\end{itemize}


\subsection{Voorkennis}
Voor een goed begrip van dit document wordt er van de lezer verwacht dat deze:
\begin{itemize}
\item kennis heeft van de werking van de shell
\item kennis heeft van de werking van commando's, opties en parameters
\item kan werken met man-pages
\end{itemize}


%%%%%%%%%%%%%%%%%
%%% De inhoud %%%
%%%%%%%%%%%%%%%%%

% Requires: fhs
% Provides: touch, cat, mkdir, rmdir, rm, redirects: >, >>, <, <<
\section{Werken met bestanden}
Als alles een bestand is in Linux dan is het werken met \index{Bestanden}bestanden het belangrijkste wat er is. Dit hoofdstuk gaat je dan ook de basisbeginselen bijbrengen van het werken met bestanden.

\subsection{Directories}
Om data op een computer te structureren is het handig om de data te verdelen over directories\index{Directories}. Directories zijn ook bekend als mappen\index{Mappen} en folders\index{Folders}. Wij zullen alleen nog spreken van directories omdat binnen de Unix-wereld de meest gebruikte term is.

Een directory maak je aan met het commando \texttt{mkdir}\index{mkdir}\index{Directories!mkdir}:
\begin{lstlisting}[language=bash]
$ mkdir LinuxCursus
\end{lstlisting}
Met \texttt{ls} kan je controleren of de directory ook daadwerkelijk aangemaakt is.

Je kan ook meerdere directories tegelijk aanmaken door ze als een lijst op te geven, gescheiden door spaties:
\begin{lstlisting}[language=bash]
$ mkdir Aap Noot Mies
\end{lstlisting}

Soms wil je ook een heel pad gelijk aanmaken met:
\begin{lstlisting}[language=bash]
$ mkdir Boom/Roos/Vis/Vuur
\end{lstlisting}
gaat dat niet lukken, want de Boom directory bestaat niet. Gelukkig kan je aan \texttt{mkdir} en optie meegeven die vertelt dat \texttt{mkdir} ook alle onderliggende directories moet aanmaken:
\begin{lstlisting}[language=bash]
$ mkdir -p Boom/Roos/Vis/Vuur
\end{lstlisting}
Gebruik \texttt{ls} om te controleren dat alle directories aanwezig zijn.

Tot slot wil je ook instaat zijn om directories weg te gooien. Met \texttt{rmdir}\index{rmdir}\index{Directories!rmdir} kan dir als de directories leeg zijn.
\begin{lstlisting}[language=bash]
$ rmdir Aap Noot Mies
\end{lstlisting}
gooit keurig alle aangemaakte directies weg. Controleer dit met \texttt{ls}. Maar doen we:
\begin{lstlisting}[language=bash]
$ rmdir Boom
\end{lstlisting}
dan krijgen we een error melding, want de \texttt{Boom} directory is niet leeg. We zullen dus eerst alle andere directories moeten weggooien, te beginnen met \texttt{Vuur}, dan \texttt{Vis}, dan \texttt{Roos} en tot slot kunnen we pas \texttt{Boom} weggooien. Gebruik \texttt{rmdir} om alle directories \textbf{behalve} Boom te verwijderen. De \texttt{Boom} directory moet dus blijven bestaan. 

Het \texttt{cd} commando gebruiken we om van directory te wisselen (cd staat voor \textbf{c}hange \textbf{d}irectory). We kunnen aan \texttt{cd} een relatief of een absoluut pad meegeven. Een absoluut pad is het complete pad vanaf de root-directory:
\begin{lstlisting}[language=bash]
$ cd /home/dennis
\end{lstlisting}
Waarbij je \textquote{dennis} vervangt door je eigen gebruikersnaam zorgt ervoor dat je in je eigen home-directory komt te staan. Als je een relatief pad gebruikt betekent dit dat je niet het hele pad meegeeft, maar een stukje. Bijvoorbeeld:
\begin{lstlisting}[language=bash]
$ cd LinuxCursus
\end{lstlisting}
Je geeft niet het complete pad \texttt{/home/dennis/LinuxCursus} op maar slechts het deel \texttt{LinuxCursus}, wat dus relatief is ten opzichte van de home-directory waarin je al staat.



\section{Bestanden maken en stdin, stdout en stderr}
Zorg dat je in de directory LinuxCursus staat en type:

\begin{lstlisting}[language=bash]
$ touch hello.txt
\end{lstlisting}

Na de Enter lijkt er helemaal niets te gebeuren. Dit is met de meeste Linux commando's het geval. Als het goed gegaan is
dan laten ze niets weten, een beetje als ``geen nieuws, is goed nieuws''. Doen we een `ls' dan zien we dat er een
bestand is aangemaakt dat hello.txt heet.

Met touch kunnen we dus bestanden aanmaken, dit zijn lege bestanden. Type maar eens:

\begin{lstlisting}[language=bash]
$ cat hello.txt
\end{lstlisting}

dan zal je zien dat er weer niets op je scherm verschijnt. En dat is goed! Het `cat' commando plaatst de inhoud van een
bestand op het scherm en daar we een leeg bestand hebben opgevraagd is wat er op het scherm komt dus niets en omdat dat
succesvol is verlopen hoeft cat ook geen foutmelding te laten zien en met de wetenschap dat geen nieuws, goed nieuws is
is cat klaar.

In de acties die we nu hebben doorlopen staat er dat echo tekst naar het scherm schrijft en cat
bestanden op het scherm afbeeldt. Dat is inderdaad wat er gebeurt, maar vanuit het Linux gezien niet helemaal correct.
Zowel echo als cat schrijven naar de standaard output en in de terminal is het scherm de standaard output. De standaard
output wordt vaak afgekort als stdout.

We kunnen de standaard output ook omleiden (redirect) naar bijvoorbeeld een bestand:

\begin{lstlisting}[language=bash]
$ echo 'Ik werk met Linux' > hello.txt
\end{lstlisting}

We zien nu dat de zin die we met echo afbeelden niet meer op het scherm verschijnt. Hij is verdwenen en er lijkt weer
helemaal niets gebeurd te zijn. Als we nu

\begin{lstlisting}[language=bash]
$ cat hello.txt
\end{lstlisting}

doen dan zien we waar onze zin is gebleven. Hij is hello.txt terecht gekomen. We hebben de stdout van echo in hello.txt
gestopt.

Laten we dat nog eens doen:

\begin{lstlisting}[language=bash]
$ echo 'Hello World!' > hello.txt
\end{lstlisting}

Doen we een cat van hello.txt dan zien we dat onze eerste zin verdwenen is en er alleen nog 'Hello World!' in hello.txt
zit. We hebben kennelijk ons bestand overschreven met nieuwe inhoud. We kunnen ook tekst toevoegen aan een bestand:

\begin{lstlisting}[language=bash]
$ cat 'Ik werk met Linux' >> hello.txt
\end{lstlisting}

door gebruik te maken van het dubbele groter dan teken voegen we een regel toe aan het eind van het bestand. De oude
regel zie je met cat als eerste en daaronder komt onze nieuwe regel.

Zou er als we een stdout hebben ook een standaard input zijn en kunnen we daar dan van lezen? Ja, die is er. Als je
typt:

\begin{lstlisting}[language=bash]
$ cat < hello.txt
\end{lstlisting}

dan vertellen we eigenlijk dat cat de invoer (stdin) op het scherm moet afbeelden. Maar ja, dan moeten we een spatie en
een <-teken extra typen en dat doen we liever niet.

Naast de standaard input en standaard output is er ook nog standaard error (stderr), waar de foutmeldingen naartoe gaan.
Laten we eens een fout maken door een niet bestaan bestand aan cat te geven:

\begin{lstlisting}[language=bash]
$ cat Hello.txt
\end{lstlisting}

Je krijgt nu een foutmelding dat Hello.txt niet bestaat. Linux is case-sensitive dus hello.txt is niet hetzelfde als
Hello.txt. De standaard output, input en error zijn genummerd in Linux. Stdin is 0, stdout is 1 en stderr is 2. Nu we
dit weten zouden we het volgende kunnen doen:

\begin{lstlisting}[language=bash]
$ cat Hello.txt 2> Hello_error.txt
\end{lstlisting}

We zien nu geen foutmelding meer op ons scherm en hebben de stderr omgeleid (redirect) naar het bestand
Hello\_error.txt. Doen we nu een
\begin{lstlisting}[language=bash]
$ cat Hello_error.txt
\end{lstlisting}
dan zien we dat de foutmelding daar is opgeslagen.

Deze vormen van het redirecten (omleiden) van data stromen gebeurt in Linux heel vaak. Programma's schrijven
bijvoorbeeld de fouten die ze tegen komen naar een log bestand. Dit doen ze door de stderr te redirecten en de error
regels toe te voegen aan het bestand, bijvoorbeeld 2>>error.log. Als ze dit doen met een
datum- en tijdmelding dan kan je heel makkelijk problemen opzoeken.

Wij hebben nu geleerd om bestanden aan te maken en om invoer en uitvoer te redirecten.


\section{Bestanden kopie\"eren, verplaatsen, hernoemen, verwijderen}
Om bestanden te kopie\"eren gebruiken \texttt{cp}\index{cp}\index{Bestanden!cp} van het Engelse copy:
\begin{lstlisting}[language=bash]
$ cp hello.txt Boom/hello.txt
\end{lstlisting}
We kunnen ook gelijk de naam veranderen als we dat willen:
\begin{lstlisting}[language=bash]
$ cp hello.txt Boom/Hallo.txt
\end{lstlisting}

Om bestanden verplaatsen gebruiken \texttt{mv}\index{mv}\index{Bestanden!mv} van het Engelse move. Het verschil met copy is dat een bestand niet meer op de oorspronkelijke plek terug te vinden is. Bij move heb je dus maar 1 bestand na de handeling, na copy heb je 2 bestanden.
\begin{lstlisting}[language=bash]
$ mv Boom/Hallo.txt .
\end{lstlisting}
Het \texttt{mv} commando kunnen we ook gebruiken om bestanden van naam te veranderen:
\begin{lstlisting}[language=bash]
$ mv Hallo.txt hallo.txt
\end{lstlisting}

Voor het weggooien van bestanden gebruiken we \texttt{rm} van remove.
\begin{lstlisting}[language=bash]
$ rm hallo.txt
\end{lstlisting}

Omdat alles een bestand is op een Linux systeem zijn ook directories bestanden, speciale bestanden, maar toch bestanden. We hebben al gezien dat we met \texttt{rmdir} lege directories weg kunnen gooien. Zouden we nu \texttt{rm} kunnen gebruiken om ook directories weg te gooien. Ja, dat kan, maar ook hier geldt dat de directory leeg moet zijn.
\begin{lstlisting}[language=bash]
$ rm Boom
\end{lstlisting}
geeft weer een foutmelding. Het systeem zegt tegen ons dat \texttt{Boom} een directory is. Als we tegen \texttt{rm} vertellen dat hij de zaken recursive weg met gooien, dan zal de hele boomstructuur wegggegooid worden:
\begin{lstlisting}[language=bash]
$ rm -r Boom
\end{lstlisting}
De \texttt{-r} optie verteld aan rm dat hij vanuit de diepste plek in de Boom directory moet beginnen met bestanden (en directories) weg te gooien net zo lang tot hij de directory Boom weg kan gooien. In ons voorbeeld is het niet veel dat \texttt{rm} te doen heeft, maar je zal niet de eerste zijn die met \texttt{-r} meer weggooid dan de bedoeling was en dat tot de ontdekking komt dat een Linux systeem geen undelete kent. Terug halen van data is er dan ook niet bij.


%%%%%%%%%%%%%%%%%%%%%
%%% Index and End %%%
%%%%%%%%%%%%%%%%%%%%%
\printindex
\end{document}

%%% Last line %%%
